%% kapa: exam=1 round=37 year=2026 period=1 form=A subject=civil_law subject_ko=민법 seq=15 number=15
%% source: Q-Net 감정평가사 기출문제 (cst00309) artlSeq=5258924 file=raw/1008/5258924/2026년 제37회 감정평가사 1차 1교시.pdf
%% conversion: kordoc 4.24.51 -> tools/gen1.py (auto); review=proofread-sonnet; flags=none
\begin{question}{15}
甲으로부터 甲 소유의 X건물을 매도하는 대리권을 수여받은 乙이 丙과 매매계약을 체결한 경우에 관한 설명으로 옳지 않은 것을 모두 고른 것은? (다툼이 있으면 판례에 따름)
\begin{framedtext}
ㄱ. 乙의 대리권은 특별한 사정이 없는 한 丙과의 계약을 해제할 권한을 포함한다. \par
ㄴ. 甲이 乙에게 대리권을 수여한 후 甲에 대하여 성년후견이 개시되면 乙의 대리권은 소멸한다. \par
ㄷ. 乙이 丙으로부터 잔금을 수령한 경우, 특별한 사정이 없는 한 乙이 잔금을 甲에게 전달하지 않았더라도 丙의 잔금지급채무는 소멸한다. \par
ㄹ. 丙이 이행불능을 이유로 매매계약을 해제한 경우, 특별한 사정이 없는 한 丙은 乙에게 채무불이행에 기한 손해배상을 청구할 수 없다. \par
\end{framedtext}
\begin{choices}
  \item ㄱ, ㄴ
  \item ㄱ, ㄷ
  \item ㄷ, ㄹ
  \item ㄱ, ㄴ, ㄹ
  \item ㄴ, ㄷ, ㄹ
\end{choices}
\end{question}
\answer{1}
